% compendium/execution-trading-costs/market-regulation.tex
% Entries migrated on 2026-10-16 from the legacy regulatory-and-structural-evolution file of the foundations
% folder (slice 35), owned by market-structure-execution-quality, one of the file's two map candidates.
% Dispositions in notes/migration-slices.md; the file verbatim in notes/artefacts/market-structure-execution-quality/.

\section{Market regulation, venue competition and circuit breakers}

\begin{qfdefinition}[Regulation as market design: who regulates, what market-design questions are, how to evaluate a rule, and the political economy of rule-making]\label{def:market-regulation-design}
\qffield{Statement} ``Regulators create and enforce rules that facilitate trading''; good rules help traders communicate, prevent fraud and make ``all things generally \dots\ as they appear'', while other rules give favoured traders, brokers or exchanges protection from competition, protect domestic or incumbent markets, express ideology or raise revenue \citep[chap.~3, sec.~3.5, pp.~59--60]{alpha-harris}. ``The stated purpose of a regulation often is not its true objective'' (p.~60). \emph{Who regulates.} Legislatures set the framework; government agencies (in the United States the SEC and CFTC, with the Federal Reserve setting speculative margins) write and enforce regulations; self-regulatory organisations (exchanges, clearinghouses and trader associations) regulate their members under rules approved by the agency; private bodies set accounting and performance standards; international organisations (IOSCO, the World Federation of Exchanges) coordinate (sec.~3.5.1--3.5.2, pp.~60--66). Harris's summary: ``Legislatures provide the regulatory framework, the SROs provide the details, and governmental regulatory agencies provide oversight'' (p.~66). \emph{Design questions} debated around 2002 included a central limit order book, quote- against order-driven systems, internalisation and preferencing, price limits and halts, floor against screen, linkages, trading hours, who sees the book, who owns market data and what may trade \citep[chap.~9, p.~202]{alpha-harris}.

\qffield{Conditions} \emph{Evaluating a rule.} Harris's normative procedure has four parts: identify all reasonable alternative policies; specify the criteria, a social welfare function with its acceptable trade-offs; evaluate each alternative by positive analysis; and choose the one with the greatest welfare; ``When policy becomes the objective rather than the means to the objective, poor results often follow'' (chap.~9, box, p.~204). His own objectives, stated as opinion, rank the private interests of utilitarian traders first, the public benefits of liquid, informative markets second, support for profit-motivated traders only when needed for those, and hostility to strategies that exploit other traders, among them manipulators, bluffers, front runners and very fast traders picking off slow dealers (sec.~9.6, pp.~216--217). The SEC's statutory direction to consider ``efficiency, competition, and capital formation'' as well as investor protection gives no guidance on trade-offs among them (chap.~3, box, p.~61). \emph{Public and private services.} Markets charge for private services such as order routing; public services such as price continuity and the regulation of insider trading, manipulation and capital structures benefit everyone, so ``unregulated competition among market centers produces fewer public services than would be socially desirable'' (chap.~26, sec.~26.5.3, p.~538).

\qffield{Practice} \emph{Political economy.} After the 1987 crash, Harris argues, regulators facing public demand for action were better off adopting circuit breakers whether or not another crash followed, and mild rather than severe ones (\emph{regulatory risk}, Tables~28-1 and~28-2); the NYSE's Rule~80A collar on index arbitrage benefited specialists, who outnumbered arbitrageurs among members (470 of 1{,}366 members were full-time specialists in November 1997, while 20 members did 95 percent of program trades), an example of rent seeking (\emph{regulatory capture}) \citep[chap.~28, sec.~28.4, pp.~578--581]{alpha-harris}. ``After a crisis, regulators often adopt regulations that have little economic value simply to respond to public demands for action'' (p.~582). \emph{Unintended consequences}: Brazil's 0.38 percent tax on financial transactions of 1997 moved institutional trading in Brazilian stocks to ADRs in New York; Bovespa volume fell from 1.2 billion reais a day in 1997 to 350 million in 2001, partly because of the 2001 crisis, and in September 2001 the government announced an exemption for stock trades (chap.~3, box, p.~66). Exposure protection is a design choice with its own trade-offs: time precedence with a meaningful tick, undisclosed orders, audit trails and delayed reporting (\hyperref[def:order-exposure-decisions]{order exposure decisions}).

\qffield{Pitfalls} The institutions named are as of 2002 (NASD has since become part of FINRA; national regulators listed in Harris's Table~3-11 have been renamed or merged). The legacy ``measurement risk'', policy adopted before transitory and fundamental volatility can be separated, matches Harris's remark that empirical studies of circuit breakers ``generally have been inconclusive'' because extreme events are rare (chap.~28, p.~577), not a named category of his. A welfare ranking is opinion; the four-step procedure only makes the criteria explicit.

\qffield{See also} \hyperref[def:market-functions]{what markets do}; \hyperref[mech:venue-competition-fragmentation]{venue competition and fragmentation}; \hyperref[mech:circuit-breakers-trading-halts]{circuit breakers and trading halts}; \hyperref[rule:us-equity-market-structure]{U.S.\ equity market structure rules}; \hyperref[rule:eu-trading-venue-transparency]{EU trading venues and transparency}.

\qffield{Further reading} Harris, chaps.~3, 9, 26 and~28. The legacy file is kept verbatim in \nolinkurl{notes/artefacts/market-structure-execution-quality/}.

\qffield{Index} market regulation; regulators; self-regulatory organisations; SEC; CFTC; IOSCO; World Federation of Exchanges; market design; welfare economics; social welfare function; normative analysis; public services; private services; regulatory risk; regulatory capture; rent seeking; unintended consequences; transaction tax

\qffield{Status} Verification: partly checked 2026-10-16 (against the text layer of a local copy of \texttt{alpha-harris}, 2002 Oxford printing, pp.~59--66, 202--218, 538 and 572--582). Stability: draft. Currency: framework timeless; institutions and examples as of 2002.
\end{qfdefinition}

\begin{qfmechanism}[Venue competition and fragmentation: the order-flow externality, why markets segment, cross-market precedence, the post-MiFID cycle of entry, convergence and consolidation, and technology fragmentation]\label{mech:venue-competition-fragmentation}
\qffield{Statement} A market is \emph{fragmented} when essentially the same thing trades in different market centres and \emph{consolidated} when all traders trade in one; competition among traders for the best price works best in a consolidated market, while competition among market centres to offer services implies fragmentation, ``The two competitions therefore are inconsistent'' \citep[chap.~26, p.~524]{alpha-harris}. \emph{Consolidating force}: each trader who joins a market adds liquidity that attracts others, the \emph{order-flow externality}, a network externality that makes trading winner-take-all and new entry hard (secs.~26.2 and~26.5.1, pp.~526 and~535--536); Hasbrouck likewise calls the liquidity externality a network externality that ``favors market consolidation'' \citep[chap.~1, sec.~1.2, p.~5]{s6-hasbrouck}. \emph{Fragmenting force}: ``Markets fragment because traders are not all identical'' (sec.~26.3, p.~530): by size (large traders control exposure and value strict time precedence with an economically significant tick; small traders want exposure and small ticks), information (informed traders prefer anonymous consolidated markets; uninformed traders prefer to be identified), patience (impatient traders value dealers; patient traders value order-driven systems with time precedence), access, and creditworthiness (a single standard makes creditworthy traders subsidise others) (secs.~26.3.1--26.3.5, pp.~530--533).

\qffield{Conditions} \emph{Order-flow externality as an inequality.} If a trader's utility at venue $i$ is $U_i=s_i-c_i+\ell(F_i)$ (service, cost, and an increasing function of order flow), an entrant better by $(s_e-c_e)-(s_m-c_m)>0$ still loses the trader when $\ell(F_m)-\ell(F_e)$ exceeds that margin; this restatement of the legacy proposition is a rearrangement, true whenever the trader chooses the higher utility, and its content lies in the specification (one trader, flows taken as given). With $\ell=\log$, flows of 1{,}000 against 50 and a 0.25 service advantage, the incumbent wins by 2.75 (\nolinkurl{slice35-checks.py}). \emph{Cross-market precedence.} A market cannot meaningfully enforce time and public-order precedence when other markets trading the same instrument do not: quote matchers trade elsewhere, small markets offer small ticks to attract them, and larger ones follow; Harris's \emph{Race to the Bottom} sidebar reads the U.S.\ tick reductions of the 1990s, which ``culminated in the full decimalization of the markets in 2001'', this way (sec.~26.5.2, pp.~536--538); his sidebar \emph{Lost Time Precedence} gives the 2002 U.S.\ example of regional specialists and third-market dealers filling orders at prices first quoted on the NYSE (p.~537). The U.S.\ answer was a trade-through rule (\hyperref[rule:us-equity-market-structure]{U.S.\ equity market structure rules}); in Europe, best-execution policies took its place \citep[chap.~1, sec.~1.1.2, p.~43]{alpha-laruelle}.

\qffield{Practice} \emph{Post-MiFID Europe, as Lehalle and Laruelle describe it.} MiFID, in force from 1 November 2007, abolished concentration rules and introduced brokers' best-execution policies; instead of many venues a few appeared, each running several pools, and ``the agents \dots\ inherited the crucial role of consolidating the liquidity'' \citep[Introduction, p.~17]{alpha-laruelle}. Their phases (chap.~1, secs.~1.1.2--1.1.4, pp.~39--62): first attempts (Chi-X, Turquoise, BATS Europe and NASDAQ OMX Europe, opened country by country, Table~1.1, p.~40); convergence towards pan-European offers and a few groups (the LSE Group with Turquoise; BATS and Chi-X; Euronext), which then segmented liquidity into niche pools, ``natural'' against ``opportunistic'' (pp.~41--42); and broker crossing networks and MTF dark pools (pp.~54--56). Their two cautionary cases: Turquoise, founded by nine banks with a market-making agreement whose end made its share drop (pp.~47--48), and NASDAQ OMX's MTF, which ``almost never crossed 1\% of the market share'' and closed (pp.~49--50). Consolidation followed: BATS acquired Chi-X Europe in November 2011, a Deutsche Börse--NYSE Euronext merger was rejected by the European Commission, ICE bought NYSE Euronext in 2013, and Cboe acquired Bats in February 2017 (pp.~60--62); their summary: challengers first offer new services at lower cost, practices converge, then mergers come while ``the offer continues to expand'' (p.~62). Harris's two consolidation waves, after the telegraph and telephone and from the early 1990s, have the same drivers: the order-flow externality, a falling ratio of variable to fixed costs, and loosened cross-border restrictions (sec.~26.6, pp.~538--539). \emph{Technology fragmentation}: application protocols (XDP, ITCH, EBS, Pitch, EMDI, RWF), multicast UDP or TCP with venue-specific recovery, and FAST, ZLIB or no compression made a common feed handler hard to build (Lehalle and Laruelle, pp.~53--54). \emph{Shared references}: in the United States the consolidated tape lets participants decide on the same information; Lehalle and Laruelle argue that Europe needs a consolidated post-trade tape rather than reliance on primary markets, citing the Euronext outage of October 2010, during which MTF prices diverged from Euronext's by about 2 percent for some stocks (Introduction, p.~7; chap.~2, sec.~2.5.2, pp.~183--186).

\qffield{Pitfalls} Lehalle and Laruelle's dates are not all consistent: the end of Turquoise's market-making agreement is March 2009 in Table~1.1 and on p.~47 but March 2008 on p.~53; NASDAQ OMX Europe runs ``from September 2008 to May 2010'' on p.~49 but launches in May 2009 and closes in July 2010 in Table~1.1; the LSE's acquisition of Turquoise is placed in February 2011 (pp.~16 and~60); and the Euronext outage is dated 13 October 2010 in the text and 10 October 2010, a Sunday, in the figure captions. These were not checked against primary sources and are reported as the book gives them. The legacy statement that segmentation separates ``natural'' from ``opportunistic'' liquidity is the authors' interpretation, with Menkveld cited for opportunistic liquidity (p.~42). Their merger-wave chronology is to 2017; current ownership was not checked. Harris's examples (ITS, SOES) are U.S.\ institutions of 2002.

\qffield{See also} \hyperref[ev:european-fragmentation-evidence]{European fragmentation evidence}; \hyperref[meth:fragmentation-concentration-metrics]{fragmentation and concentration metrics}; \hyperref[def:smart-order-routing]{smart order routing}; \hyperref[def:dark-pools-crossing]{dark pools and crossing networks}; \hyperref[def:market-structures-venues]{market structures and venues}; \hyperref[def:fixed-income-market-structure]{fixed-income market structure}.

\qffield{Further reading} Harris, chap.~26; Lehalle and Laruelle, introduction and chap.~1, sec.~1.1. The legacy file is kept verbatim in \nolinkurl{notes/artefacts/market-structure-execution-quality/}.

\qffield{Index} fragmentation; consolidation; segmentation; order-flow externality; network externality; winner-take-all; secondary precedence rules; time precedence across markets; quote matching; race to the bottom; MiFID; MTF; Chi-X; Turquoise; BATS Europe; NASDAQ OMX Europe; broker crossing networks; mergers of exchanges; feed protocols; consolidated tape

\qffield{Status} Verification: partly checked 2026-10-16 (against the text layers of local copies of \texttt{alpha-harris}, 2002 Oxford printing, pp.~524--541, \texttt{s6-hasbrouck}, 2007 Oxford printing, p.~5, and \texttt{alpha-laruelle}, 2018 World Scientific printing, pp.~4--7, 15--17, 39--62 and 177--186). Stability: draft. Currency: mechanisms timeless; European chronology to 2017, as reported.
\end{qfmechanism}

\begin{qfrule}[U.S.\ equity market structure rules: decimal pricing, Regulation NMS and the minimum pricing increment, single-stock and market-wide circuit breakers, stub quotes, market access controls and the consolidated audit trail (status as of 16 October 2026)]\label{rule:us-equity-market-structure}
\qffield{Statement} \emph{Decimal pricing.} On 8 June 2000 the SEC ordered the exchanges and the NASD to submit a plan phasing in decimal pricing for stocks and options, starting no later than 5 September 2000 \citep[Release 34-42914]{sec-34-42914}; the conversion of consolidated quotations ``was successfully completed on April 9, 2001'', narrowing the minimum price variation from $1/16$ of a dollar to a penny \citep[Release 34-44568, p.~38390]{sec-34-44568}. \emph{Regulation NMS} (Release 34-51808, effective 29 August 2005) adopted the Order Protection Rule (policies and procedures to prevent trades at prices inferior to protected quotations, which must be immediately and automatically accessible), the Access Rule (fair access, a cap on access fees, and a ban on patterns of locking or crossing quotations), the Sub-Penny Rule (no orders, quotes or indications in increments below a penny except when priced below \$1.00) and amended market data rules \citep[p.~37496]{sec-34-51808}; its phased compliance dates were not read. \emph{Minimum pricing increment, current text.} As amended in 2024 \citep[Release 34-101070, catalogue entry]{ex-sec-tick-access-2024}, 17 CFR 242.612 sets, for orders priced at \$1.00 or more, an increment of \$0.01 if a stock's time-weighted average quoted spread over the evaluation period (January--March or July--September) exceeded \$0.015 and \$0.005 if it did not, operative from the first business day of May and of November; \$0.0001 below \$1.00 \citep[text for 14 October 2026]{cfr-242-612}.

\qffield{Conditions} \emph{Adoption, stay and compliance of the 2024 amendments.} Adopted 18 September 2024 (Release 34-101070), the amendments also cut the access-fee caps of Rule 610(c) to \$0.001 per share for quotations at \$1.00 or more and 0.1 percent of the price below, required exchange fees and rebates to be determinable at execution (Rule 610(d)), and accelerated the round-lot and odd-lot information definitions. On 12 December 2024 the Commission stayed the amendments to Rules 600(b)(89)(i)(F), 610(c) and 612 pending judicial review, not those to 610(d), 603(b) or the round-lot and odd-lot definitions; the D.C.\ Circuit denied the petitions on 14 October 2025; on 31 October 2025 the Commission exempted compliance with the three stayed rules until the first business day of November 2026; and on 11 June 2026 it extended the exemption to ``the first business day of November 2027'' for the half-penny increment and the lower fee caps \citep[pp.~1--5]{sec-34-105656}. As of 16 October 2026 the amended Rule 612 is therefore in the CFR but not yet operative; the order cites other initiatives scheduled for 2026 and a staff review of Rules 610(c) and 612. \emph{Circuit breakers and quoting after 6 May 2010.} Single-stock circuit breakers (five-minute pauses after a 10 percent move in five minutes) were adopted in June 2010 for S\&P 500 stocks and extended to all NMS securities in June 2011, per Lehalle and Laruelle \citep[chap.~2, sec.~2.5.2, pp.~181--182]{alpha-laruelle}. On 31 May 2012 the Commission approved, as a one-year pilot, the NMS Plan to Address Extraordinary Market Volatility (\emph{limit up--limit down}), which prevents trades outside price bands around a reference price, with bands, for stocks priced above \$3.00, of 5 percent for Tier~1 stocks (S\&P 500, Russell 1000, certain ETPs) and 10 percent for Tier~2 stocks, 20 percent from \$0.75 to \$3.00, and the lesser of \$0.15 or 75 percent below \$0.75, with a proposed implementation date of 4 February 2013 \citep[Release 34-67091]{sec-34-67091}; on 11 April 2019 it approved the plan's operation ``on a permanent basis'' \citep[Release 34-85623]{sec-34-85623}. On the same day in 2012 it approved, as a pilot, revised \emph{market-wide circuit breakers}: the S\&P 500 replaced the Dow Jones Industrial Average, triggers recalculated daily fell from 10, 20 and 30 percent to 7, 13 and 20 percent, Level~1 and~2 declines halt trading for 15 minutes before 3:25 p.m.\ and not after, and Level~3 halts trading for the day \citep[Release 34-67090]{sec-34-67090}; exchange filings extended that pilot in 2019 and 2020 (to 18 October 2021 in one filing) \citep[MEMX filing, sec.~II.A.1]{sec-fr-2020-23011}, and its current status was not checked. \emph{Stub quotes.} The Commission approved SRO rules, implemented 6 December 2010, requiring registered market makers to quote within 8 percent of the NBBO for securities under the single-stock circuit breakers (readjusting at 9.5 percent), 20 percent outside the breakers' hours and 30 percent for other securities \citep[Release 34-63255]{sec-34-63255}. \emph{Market access.} Rule 15c3-5 (Release 34-63241, dated 3 November 2010; effective 14 January 2011, compliance 14 July 2011) requires brokers with market access to control pre-set credit and capital thresholds and regulatory requirements under their own direct control \citep[paras.~(b)--(c)]{cfr-240-15c3-5}. \emph{Consolidated audit trail.} Rule 613 (Release 34-67457, effective 1 October 2012) required the SROs to submit an NMS plan for a consolidated order-tracking system from order inception through routing, cancellation, modification or execution \citep[p.~45722, summary and dates]{sec-34-67457}; the CAT NMS Plan was approved on 15 November 2016 \citep[Release 34-79318]{sec-34-79318}; its current amendments and cost litigation \citep[catalogue entry only]{wide-sec-cat} were not read.

\qffield{Practice} Historical claims to keep apart: Harris's account of trade-through and time-precedence problems before Regulation NMS (\hyperref[mech:venue-competition-fragmentation]{venue competition and fragmentation}) and his description of market-wide halts at 10, 20 and 30 percent of the DJIA after 1998 \citep[chap.~28, box, p.~573]{alpha-harris} predate the 2012 rules. Lehalle and Laruelle date the stub-quote approval 8 November 2010; the approval order in the Federal Register is dated 5 November 2010. They date the vote on Rule 613 11 July 2012; the release's own date was not confirmed. The legacy marketrule ``Decimalization \dots\ narrowed quoted spreads but reduced displayed depth and queue stability'' is evidence, summarised in \hyperref[ev:tick-size-market-quality]{tick size and market quality}; ``Reg NMS links U.S.\ equity venues through protected quotes'' omits that protection covers only automated top-of-book quotations (Rule 611, \hyperref[def:best-execution]{best execution}).

\qffield{Pitfalls} The compliance status of the 2024 tick and access-fee amendments has changed three times; recheck the SEC's exemptive orders before relying on a date. Lehalle and Laruelle's single-stock circuit-breaker parameters were not checked against the 2010 approval orders. Rule texts are from the eCFR point-in-time edition for 14 October 2026; SEC releases were read in the Federal Register or, for the 2026 order, in the SEC's PDF.

\qffield{See also} \hyperref[def:best-execution]{best execution}; \hyperref[def:tick-size]{tick size}; \hyperref[mech:circuit-breakers-trading-halts]{circuit breakers and trading halts}; \hyperref[def:hft-strategies]{high-frequency trading strategies}; \hyperref[rule:eu-trading-venue-transparency]{EU trading venues and transparency}.

\qffield{Further reading} Lehalle and Laruelle, chap.~2, sec.~2.5.2. The legacy file is kept verbatim in \nolinkurl{notes/artefacts/market-structure-execution-quality/}.

\qffield{Index} decimalisation; decimal pricing; minimum price variation; Regulation NMS; Order Protection Rule; trade-through; Access Rule; access fee cap; Sub-Penny Rule; Rule 612; half-penny tick; time-weighted average quoted spread; exemptive relief; limit up--limit down; market-wide circuit breakers; single-stock circuit breakers; stub quotes; market maker quoting obligations; Rule 15c3-5; market access; Rule 613; consolidated audit trail

\qffield{Status} Verification: partly checked 2026-10-16 (Federal Register texts of Releases 34-42914, 34-44568, 34-51808 (first pages and dates), 34-63255, 34-67090, 34-67091, 34-67457, 34-79318 (date) and 34-85623 (date and purpose), and of an exchange filing of October 2020; SEC Release 34-105656 of 11 June 2026, pp.~1--5; 17 CFR 242.612 and 240.15c3-5 in the eCFR for 14 October 2026; \texttt{alpha-laruelle}, pp.~181--183; \texttt{alpha-harris}, p.~573). Stability: draft. Currency: as of 16 October 2026; market-wide circuit-breaker and CAT status after 2021 unchecked.
\end{qfrule}

\begin{qfrule}[EU trading venues and pre-trade transparency: MiFID and MiFID II dates, venue types, waivers, and the double volume cap replaced by a single volume cap (status as of 16 October 2026)]\label{rule:eu-trading-venue-transparency}
\qffield{Statement} \emph{Dates.} MiFID (Directive 2004/39/EC of 21 April 2004) applied from 1 November 2007, after Directive 2006/31/EC postponed its transposition deadline to 31 January 2007 \citep[Art.~1, points (4)--(5)]{eu-dir-2006-31}. MiFID~II (Directive 2014/65/EU) and MiFIR (Regulation (EU) No 600/2014) were to apply from 3 January 2017 \citep[Art.~93]{eu-dir-2014-65} \citep[Art.~55]{eu-reg-600-2014}; Directive (EU) 2016/1034 and Regulation (EU) 2016/1033, both of 23 June 2016, deferred them by twelve months to 3 January 2018 \citep[Art.~1, points (7)--(8)]{eu-dir-2016-1034} \citep[recital and Art.~1]{eu-reg-2016-1033}. \emph{Venue types} under MiFID~II: a \emph{regulated market} and an \emph{MTF} are multilateral systems bringing together third-party interests under non-discretionary rules; an \emph{OTF} is a multilateral system that is neither, for bonds, structured finance products, emission allowances or derivatives; a \emph{trading venue} is any of the three; a \emph{systematic internaliser} is an investment firm dealing on own account when executing client orders outside them ``on an organised, frequent systematic and substantial basis'' without operating a multilateral system \citep[Art.~4(1), points (20)--(24)]{eu-dir-2014-65}. \emph{Dark trading} uses the \emph{reference price waiver} and the \emph{negotiated trade waiver} of MiFIR Article 4(1)(a) and (b)(i) \citep[recital on the volume cap]{ex-eu-mifir-review-2024}.

\qffield{Conditions} \emph{From double to single volume cap.} Under the original MiFIR Article 5, trading under those two waivers in an instrument was limited to 4 percent of total EU trading on any one venue and 8 percent across the EU over the previous 12 months; breaching either led the competent authorities to suspend the waivers within two working days for six months, on monthly ESMA data, with warning reports at 3.75 and 7.75 percent \citep[Art.~5]{eu-reg-600-2014}. Regulation (EU) 2024/791 of 28 February 2024, published on 8 March 2024 and in force from 28 March 2024, replaced it with a single cap: trading venues suspend their own use of the reference price waiver for three months when trading under it exceeds 7 percent of total EU trading in the instrument, based on quarterly ESMA data covering the previous 12 months; the negotiated trade waiver is no longer capped; monitoring under the new cap ``shall start on 29 September 2025''; ESMA reports annually from 29 September 2027 and the Commission may adjust the threshold \citep[Art.~1(4), amending MiFIR Art.~5]{ex-eu-mifir-review-2024}. \emph{Transition.} ESMA stated on 27 March 2024 that the double cap would continue, for both waivers, until the single cap applied \citep[p.~2]{esma-transition-statement-2024}; its double-volume-cap page records that the double cap was discontinued on 28 September 2025, ending any ongoing suspension, and that the first single-cap file was planned for 9 October 2025 \citep[retrieved 2026-10-16]{esma-dvc-page}; on 10 October 2025 ESMA confirmed that the single cap ``has become operational with the first calculation results published on 9 October 2025'' \citep[p.~3]{esma-mifir-review-statement-2025}.

\qffield{Practice} \emph{Lehalle and Laruelle's account (2018).} MiFID's ``few and simple mechanisms'' were best execution, three venue statuses and transparency (chap.~1, sec.~1.1.2, p.~42): regulated markets with public books, responsible for circuit breakers; MTFs, often pegged to regulated markets; and systematic internalisers, which few firms chose because of the transparency attached (pp.~43--44); the pre-trade waivers were large in scale, order management facility, negotiated trade and reference price, under which dark pools operated (p.~45). Writing before MiFID~II applied, they expected a dark cap of 4 and 8 percent enforced yearly with six-month suspensions and a tick-size regime set by price and trading activity (Introduction, pp.~8--9); the first matched the original Article~5, the second is MiFID~II Article~49, which requires tick-size regimes calibrated to liquidity and average spread \citep[Art.~49]{eu-dir-2014-65}. \emph{Payment for order flow}: the 2024 amendments prohibit investment firms from receiving payment for forwarding or executing retail and opted-in professional client orders, with a Member-State exemption possible until 30 June 2026 for firms already doing so \citep[Art.~1, new MiFIR Art.~39a]{ex-eu-mifir-review-2024}.

\qffield{Pitfalls} Do not quote ``4 and 8 percent'' as current: they applied from January 2018 to 28 September 2025. The single cap is 7 percent, applies only to the reference price waiver, is triggered by venues themselves, and suspends for three months. ESMA's data and suspension files are the operational record; individual suspensions were not checked. The U.K.\ regime, which diverged after 2020, was not checked. Article~48 of MiFID~II (venue circuit breakers) is in \hyperref[mech:circuit-breakers-trading-halts]{circuit breakers and trading halts}; best execution is \hyperref[def:best-execution]{best execution}.

\qffield{See also} \hyperref[def:dark-pools-crossing]{dark pools and crossing networks}; \hyperref[rule:us-eu-bond-transparency]{bond transparency rules}; \hyperref[mech:venue-competition-fragmentation]{venue competition and fragmentation}; \hyperref[def:tick-size]{tick size}; \hyperref[rule:us-equity-market-structure]{U.S.\ equity market structure rules}.

\qffield{Further reading} Lehalle and Laruelle, introduction and chap.~1, sec.~1.1. The legacy file is kept verbatim in \nolinkurl{notes/artefacts/market-structure-execution-quality/}.

\qffield{Index} MiFID; MiFID II; MiFIR; MiFIR review; Regulation (EU) 2024/791; regulated market; multilateral trading facility; organised trading facility; systematic internaliser; trading venue; pre-trade transparency waivers; reference price waiver; negotiated trade waiver; large in scale; double volume cap; single volume cap; ESMA; tick-size regime; payment for order flow

\qffield{Status} Verification: partly checked 2026-10-16 (texts of Directives 2006/31/EC, 2014/65/EU and (EU) 2016/1034 and Regulations (EU) No 600/2014, 2016/1033 and 2024/791, retrieved through the EU Publications Office's CELLAR service because EUR-Lex's HTML pages returned an empty bot-check response; ESMA's statements of 27 March 2024 and 10 October 2025 and its double-volume-cap page; \texttt{alpha-laruelle}, pp.~8--9 and 42--45). Stability: draft. Currency: as of 16 October 2026.
\end{qfrule}

\begin{qfmechanism}[Circuit breakers and trading halts: halts, price limits, taxes, margins and collars; when a halt helps; the gravitational effect, propagation and related instruments; incidence by trader type]\label{mech:circuit-breakers-trading-halts}
\qffield{Statement} Harris: ``Circuit breakers are trading rules that limit trading activity'': trading halts, price limits, transaction taxes, margin requirements and position limits, and collars that restrict access to trading systems \citep[chap.~28, sec.~28.3, p.~572]{alpha-harris}. Lehalle and Laruelle classify thresholds as \emph{static} (``a fixed price computed before the opening'') or \emph{dynamic} (``a corridor following the recent values of the price''), and breakers as \emph{local} (one instrument) or \emph{global} (a collection, generally all) \citep[chap.~2, sec.~2.5.3, pp.~187--188]{alpha-laruelle}. MiFID~II requires regulated markets to be able ``to temporarily halt or constrain trading if there is a significant price movement in a financial instrument on that market or a related market during a short period'', with calibrated parameters reported to ESMA, and a material venue that halts must let authorities coordinate a market-wide response \citep[Art.~48(5)]{eu-dir-2014-65}. U.S.\ rules are in \hyperref[rule:us-equity-market-structure]{U.S.\ equity market structure rules}.

\qffield{Conditions} \emph{When a halt helps} (Harris, sec.~28.3.1, pp.~572--575). If new fundamental information drives the move, a halt ``will merely postpone the inevitable'' and makes prices less informative meanwhile; if an imbalance among uninformed traders drives it, stopping them may help. Halts can let liquidity suppliers arrive, publicise imbalances, give time to meet margin calls and remove stop orders, and switch the pricing rule from discriminatory (continuous trading) to uniform (the reopening call), protecting standing limit orders. Against them: the \emph{gravitational effect}, traders who fear a halt ``may submit their orders earlier'', so a price-triggered rule can raise volatility near its trigger, which discretionary, less predictable halts avoid; and value traders may watch less if they expect halts. A halt in one market diverts order flow to others trading the same risk, so uncoordinated breakers ``may increase volatility''. Harris's conclusion: the theoretical net effect on transitory volatility is indeterminate, and empirical studies ``generally have been inconclusive'' because extreme events are rare (sec.~28.3.4, pp.~576--577). The legacy proposition that a halt improves price quality iff its benefit $B$ exceeds its cost $C$ is a definition of improvement, not a result; its content is the list of benefits and costs above.

\qffield{Practice} \emph{Price limits and margin.} Harris's sidebar: a trader with \$100{,}000 margin and \$100{,}000 other assets faces a 25-point fall worth \$250{,}000; without limits the broker immediately loses the margin and must chase \$150{,}000; with a five-point daily limit the trader, unsure whether he is bankrupt, may post margin for two days before defaulting, and owes \$50{,}000 at the end (box ``Am I Bankrupt?'', p.~574); ``Price limits can protect the clearing system by increasing the margin payments that bankrupt traders pay before they default'' (p.~582), and markets with halts ``can have lower margin requirements'' (p.~575). \emph{Incidence of taxes and margins} (sec.~28.3.2, pp.~575--576): restricting informed traders makes prices less informative and restricting value traders, dealers and arbitrageurs raises transitory volatility, while restricting order anticipators, bluffers and gamblers lowers it; restricting utilitarian traders would be ``highly undesirable''; and it is hard to design rules that discriminate among these types, a tax falling hardest on frequent traders, dealers included. \emph{Design}: Lehalle and Laruelle ask that breakers respond to liquidity imbalance, although ``all of them only look at prices'', perhaps by combining large moves with poor liquidity; that thresholds reflect the volatility level and its intraday seasonality; and that instruments indexed on halted assets, highly correlated instruments and other venues trading the same instrument be halted in step, which they leave as ``a research question''; they cite China's removal of badly calibrated thresholds in January 2016 (pp.~188--189). \emph{Propagation}: in the Flash Crash of 6 May 2010 the absence of a stock-specific breaker was named among aggravating factors (p.~181).

\qffield{Pitfalls} Harris's examples (the DJIA 10/20/30 percent halts, NYSE Rule 80A collars) are rules of 1998--2002 that have since been replaced (\hyperref[rule:us-equity-market-structure]{U.S.\ rules}). ``Collars'' in Harris restrict access to trading systems, as Rule 80A did for index-arbitrage program trades; the legacy wording ``restrict specific order types or program trades'' fits that case. The legacy design principle that breakers ``should stop disorderly liquidity demand, not ordinary price discovery'' restates Harris's fundamental-versus-transitory distinction; neither source shows how to tell the two apart in real time. The legacy incidence table assigned ``gamblers: discouraging stale-information trading can reduce noise''; Harris's point is that gamblers usually trade on stale information, so discouraging them would reduce transitory volatility (p.~576).

\qffield{See also} \hyperref[mech:bubble-crash-dynamics]{bubbles and crashes}; \hyperref[def:volatility-components]{fundamental and transitory volatility}; \hyperref[def:market-regulation-design]{regulation as market design}; \hyperref[mech:futures-margin-settlement]{futures margin and settlement}; \hyperref[def:trading-sessions-execution-systems]{trading sessions and execution systems}.

\qffield{Further reading} Harris, chap.~28; Lehalle and Laruelle, chap.~2, sec.~2.5. The legacy file is kept verbatim in \nolinkurl{notes/artefacts/market-structure-execution-quality/}.

\qffield{Index} circuit breakers; trading halts; price limits; limit moves; transaction taxes; margin requirements; position limits; collars; static thresholds; dynamic thresholds; local and global breakers; gravitational effect; magnet effect; halt coordination; related instruments; reopening call; uniform pricing; margin collection; incidence of trading restrictions; MiFID II Article 48

\qffield{Status} Verification: partly checked 2026-10-16 (against the text layers of local copies of \texttt{alpha-harris}, 2002 Oxford printing, pp.~572--582, and \texttt{alpha-laruelle}, 2018 World Scientific printing, pp.~181 and 187--189; Directive 2014/65/EU, Art.~48(5), via the Publications Office). Stability: draft. Currency: mechanisms timeless; examples dated as stated.
\end{qfmechanism}
